\thispagestyle{empty}
\vspace*{\fill}

{\footnotesize
\noindent
\textbf{Linux forensics: a real case}\\
\textit{Autopsy of a Cyberattack on Live Production Servers}\\
\textit{Incident response, volatile memory, and multi-system event correlation}

\vspace{1.0em}
\noindent
\textbf{Second Edition} (Open digital edition, published without commercial ISBN)

\vspace{0.8em}
\noindent
\textbf{Author:} Alexis Real Santillán\\
\textbf{Professional Profile:} Cybersecurity, Cloud Architecture, and Critical Infrastructure Consultant.\\
\textbf{Technical Leadership:} CTO at Linux Open Source de México (\href{https://linuxopensource.org}{linuxopensource.org})\\
\textbf{Curriculum Vitae and Portfolio:} \href{https://alexisreal.github.io/cv/}{alexisreal.github.io/cv/}\\
\textbf{Contact:} \href{mailto:alexisreal@gmail.com}{alexisreal@gmail.com}

\vspace{1.0em}
\noindent
\textbf{Edition History:}\\
\textbf{First Edition (institutional forensic report):} September 2005\\
\textbf{Second Edition (revised and expanded):} 2025

\vspace{1.0em}
\noindent
\textbf{Origin of the Text:}\\
This work is an educational adaptation and technical expansion of the official forensic examination report submitted to institutional authorities by the author during his tenure as Computer Systems Administrator at the Institute of Physics and Mathematics (IFM), Universidad Michoacana de San Nicolás de Hidalgo (UMSNH), an institution where he no longer works and maintains ties exclusively as an external technical consultant.

\vspace{1.0em}
\noindent
\textbf{Legal Notice and License:}\\
This work is published under an international \textbf{Creative Commons Attribution-ShareAlike 4.0 International License (CC BY-SA 4.0)}. You are free to share and adapt the material under appropriate attribution terms and the same license.

\vspace{1.0em}
\noindent
\textbf{Editorial Credits and Technical Assistance:}\\
The original edition of this work was conceived and written in Spanish. This English edition was translated with the assistance of Google Gemini. All illustrations, technical architecture diagrams, and typographical style revisions were produced with the assistance of Google Gemini.

\vspace{1.0em}
\noindent
\textbf{Editorial Note on Anonymization and Privacy:}\\
This book documents an investigation conducted on real-world production systems. In accordance with legal, technical, and ethical standards, strict anonymization guidelines have been applied:
\begin{itemize}\setlength{\itemsep}{0pt}\setlength{\parskip}{0pt}
    \item IP addresses have been mapped to documentation address spaces under \textbf{RFC 5737} (\texttt{192.0.2.0/24}, TEST-NET-1) and \textbf{RFC 2544}, strictly preserving subnet relationships and forensic correlation.
    \item Domain and machine names have been normalized under the standard academic naming convention (\texttt{ifm.lab.edu.mx}), adopting identifiers inspired by servers of that era (\texttt{helios}, \texttt{atlas}, \texttt{io}), in compliance with \textbf{RFC 2606}.
    \item Individual names and user accounts have been pseudonymized to protect personal data confidentiality.
\end{itemize}

\vspace{0.8em}
\noindent
\textit{Typeset in \LaTeX{} using TeX Gyre Pagella, TeX Gyre Heros, and Inconsolata fonts.}
}
\newpage
